\label{lecture03}
\subsection*{Цели и маршрут занятия (90 минут)}
Лекция продолжает обзор IEEE 754 из занятия №2
(с.~\pageref{lecture02-float}). Теперь важно понять не только, \emph{как
хранится} число, но и \emph{как получается} результат операции.
После занятия студент сможет проследить выравнивание порядков,
нормализацию и округление, объяснить поглощение и неассоциативность,
отличить ошибку исходных данных от ошибки операции и выбрать допуск
для сравнения приближений.

\begin{center}
\begin{tabular}{|l|r|}
\hline
Блок & Минуты от начала \\
\hline
Мотивация: неожиданная сумма & 0--5 \\
Повторение формата и точности & 5--15 \\
Сложение и округление & 15--35 \\
Вычитание и ошибки приближений & 35--50 \\
Умножение и деление & 50--62 \\
Последствия и границы диапазона & 62--78 \\
Сравнение с допуском & 78--87 \\
Итоги и самопроверка & 87--90 \\
\hline
\end{tabular}
\end{center}
Подробные битовые раскладки и расширенные случаи в конце лекции
относятся к дополнительному чтению и не входят в 90 минут.

\subsection{Мотивация: почему сумма отличается от ожидания}
В обычной реализации C с \texttt{double} формата binary64, округлением
к ближайшему и вычислениями именно в этом формате:
\begin{lstlisting}[language=C,basicstyle=\ttfamily\small]
double a = 0.1, b = 0.2, c = 0.3;
printf("%.17g\n", a + b);       // 0.30000000000000004
printf("%d\n", a + b == c);    // 0
\end{lstlisting}
Для запуска нужны \texttt{\#include <stdio.h>} и функция \texttt{main}.
Здесь происходит не одно, а несколько округлений: сначала десятичные
литералы превращаются в двоичные приближения, затем округляется сумма.
Отдельно округлённое \texttt{0.3} не обязано совпасть с суммой
приближений \texttt{0.1} и \texttt{0.2}.

Это не свойство названия языка. Например, в binary32 при тех же правилах
\texttt{0.1f + 0.2f == 0.3f} даёт истину: округления случайно приводят
к одному значению. Это не делает десятые доли точно представимыми.
\textbf{Всегда называем формат и режим вычисления.}

\subsection{Опорная справка: формат, точность и шаг}
Нормальное двоичное число имеет вид
\[
 x=(-1)^s M2^E,\qquad M=1.f_2,\qquad 1\le M<2.
\]
Значащую часть $M$ также называют мантиссой (significand); она включает
скрытую единицу, а поле $f$ содержит только дробные биты.

\begin{center}
\begin{tabular}{|l|r|r|r|}
\hline
Формат & Биты $f$ & Точность $p$ & Смещение \\
\hline
\texttt{float32} / binary32 & 23 & 24 & 127 \\
\texttt{double} / binary64 & 52 & 53 & 1023 \\
\hline
\end{tabular}
\end{center}
Для normal истинный порядок равен $E=e-\mathrm{bias}$, где $e$ ---
хранимое поле. В binary32 $1\le e\le254$, то есть $-126\le E\le127$.
При $e=0$ скрытой единицы нет: $x=(-1)^s(0.f_2)2^{-126}$.
Нулевая дробь даёт знаковый ноль, ненулевая --- субнормальное число.
При $e=255$ нулевая дробь означает бесконечность, ненулевая --- NaN.
Кодирование $-12.625$ уже разобрано в лекции №2; здесь используем
тот же способ чтения полей, не повторяя весь алгоритм.

\paragraph{Числа на неравномерной сетке.}
В нормальном диапазоне $[2^E,2^{E+1})$ расстояние между соседними
положительными значениями равно $2^{E-(p-1)}$. Например, в binary32:
\[
\begin{array}{c|c}
\text{Диапазон} & \text{Шаг} \\
\hline
{[1,2)} & 2^{-23} \\
{[2,4)} & 2^{-22} \\
{[2^{24},2^{25})} & 2
\end{array}
\]
В точке степени двойки шаг слева может отличаться от шага справа.
Машинный эпсилон $\varepsilon=2^{-(p-1)}$ --- шаг \emph{справа от 1}.
Он не равен минимальному положительному числу и не задаёт постоянную
абсолютную ошибку на всей числовой оси.

\subsection{Модель вычисления: точный результат, затем округление}
Обозначим через $\operatorname{RN}$ округление к ближайшему
представимому значению; при равных расстояниях выбираем значение
с чётным младшим сохраняемым разрядом (ties to even).
Для базовых операций IEEE 754 полезна модель
\[
 \operatorname{fl}(a\circ b)=\operatorname{RN}(a\circ b),
 \qquad \circ\in\{+,-,\times,/\}.
\]
Справа действие математически точное, а $a,b$ --- \emph{уже представимые}
операнды. Процессор не обязан физически вычислять бесконечную дробь:
он должен получить результат, эквивалентный такому округлению.
Особые значения обрабатываются по отдельным правилам.

В ручных коротких примерах будем явно указывать \textbf{учебную точность
$p=4$}: одна ведущая единица и три дробных бита. Это не binary32;
порядок в таких примерах не ограничиваем. Дополнительные биты
промежуточного результата не входят в точность хранимого числа.

Для программных примеров предполагаются binary32/binary64, режим RN,
округление после указанных операций, постепенный переход к нулю и
отключённые преобразования \texttt{fast-math}. Язык C сам по себе
не гарантирует все эти свойства: их задают реализация и параметры
компиляции. В x86 вычисления \texttt{float}/\texttt{double} обычно
выполняются инструкциями SSE, например \texttt{ADDSS}/\texttt{ADDSD}.
